% !TEX spellcheck = en_US
% !TEX spellcheck = LaTeX
\documentclass[letterpaper,10pt,english]{article}
\usepackage{%
amsfonts,%
amsmath,%
amssymb,%
amsthm,%
babel,%
bbm,%
%biblatex,%
caption,%
centernot,%
color,%
enumerate,%
%enumitem,%
epsfig,%
epstopdf,%
etex,%
fancybox,%
framed,%
fullpage,%
%geometry,%
graphicx,%
hyperref,%
latexsym,%
mathptmx,%
mathtools,%
multicol,%
paralist,%
pgf,%
pgfplots,%
pgfplotstable,%
pgfpages,%
proof,%
psfrag,%
%subfigure,%
tcolorbox,%
tfrupee,%
tikz,%
times,%
%ulem,%
url,%
xcolor,%
mathpazo,%
}

\definecolor{shadecolor}{gray}{.95}%{rgb}{1,0,0}
\usepackage[margin=1in,top=0.75in]{geometry}
\usepackage[mathscr]{eucal}
\usepgflibrary{shapes}
\usepgfplotslibrary{fillbetween}
\usetikzlibrary{%
arrows,%
backgrounds,%
chains,%
decorations.pathmorphing,% /pgf/decoration/random steps | erste Graphik
decorations.text,% 
matrix,%
positioning,% wg. " of "
fit,%
patterns,%
petri,%
plotmarks,%
scopes,%
shadows,%
shapes.misc,% wg. rounded rectangle
shapes.arrows,%
shapes.callouts,%
shapes%
}

%\pgfplotsset{compat=newest} %<------ Here
\pgfplotsset{compat=1.11} %<------ Or use this one

\theoremstyle{plain}
\newtheorem{thm}{Theorem}[section]
\newtheorem{lem}[thm]{Lemma}
\newtheorem{prop}[thm]{Proposition}
\newtheorem{cor}[thm]{Corollary}
\newtheorem{clm}[thm]{Claim}

\theoremstyle{definition}
\newtheorem{defn}[thm]{Definition}
\newtheorem{conj}[thm]{Conjecture}
\newtheorem{exmp}[thm]{Example}
\newtheorem{axiom}[thm]{Axiom}
\newtheorem{assum}[thm]{Assumption}
\newtheorem{exerc}[thm]{Exercise}
\newtheorem*{defin}{Definition}

\theoremstyle{remark}
\newtheorem{rem}{Remark}
\newtheorem{note}{Note}

\newcommand{\iid}{\emph{i.i.d.}\ }
\newcommand{\Cov}{\operatorname{Cov}}
%\newcommand{\det}{\operatorname{det}}
%\newcommand{\Real}{\mathbb{R}}
\newcommand{\tr}{\operatorname{tr}}
\newcommand{\Var}{\operatorname{Var}}
\newcommand{\cov}{\operatorname{cov}}

%\renewcommand{\proof}[1]{\begin{proof}#1\end{proof}}
\newcommand{\EQ}[1]{\begin{equation*}#1\end{equation*}}
\newcommand{\EQN}[1]{\begin{equation}#1\end{equation}}
\newcommand{\eq}[1]{\begin{align*}#1\end{align*}}
\newcommand{\meq}[2]{\begin{xalignat*}{#1}#2\end{xalignat*}}
\newcommand{\norm}[1]{\left\lVert#1\right\rVert}
\newcommand{\inner}[1]{\left\langle #1\right\rangle}
\newcommand{\abs}[1]{\left\lvert#1\right\rvert}
\newcommand{\expect}[1]{\mathbb{E}\left[{#1}\right]}
\newcommand{\prob}[1]{\mathbb{P}\left[{#1}\right]}
\newcommand{\given}{\; \big\vert \;} 
\newcommand{\set}[1]{\left\{#1\right\}} 
\newcommand{\indicator}[1]{1_{\set{#1}}} 
\newcommand{\Ind}[1]{\mathbbm{1}_{#1}} 
\newcommand{\SetIn}[1]{\mathbbm{1}_{\set{#1}}} 
\newcommand{\red}[1]{\textcolor{red}{#1}} 
\newcommand{\floor}[1]{\left\lfloor#1\right\rfloor}
\newcommand{\ceil}[1]{\left\lceil#1\right\rceil}


\newcommand{\C}{\mathbb{C}}
\newcommand{\D}{\mathbb{D}}
\newcommand{\E}{\mathbb{E}}
\newcommand{\F}{\mathbb{F}}
\newcommand{\N}{\mathbb{N}}
\renewcommand{\P}{\mathbb{P}}
\newcommand{\Q}{\mathbb{Q}}
\newcommand{\R}{\mathbb{R}}
\newcommand{\Z}{\mathbb{Z}}

\newcommand{\bA}{\mathbf{A}}
\newcommand{\bI}{\mathbf{I}}
\newcommand{\bU}{\mathbf{1}}
\newcommand{\bx}{\mathbf{x}}
\newcommand{\bX}{\mathbf{X}}
\newcommand{\bZ}{\mathbf{Z}}


\newcommand{\cA}{\mathcal{A}}
\newcommand{\cB}{\mathcal{B}}
\newcommand{\cC}{\mathcal{C}}
\newcommand{\cF}{\mathcal{F}}
\newcommand{\cG}{\mathcal{G}}
\newcommand{\cM}{\mathcal{M}}
\newcommand{\cN}{\mathcal{N}}
\newcommand{\cP}{\mathcal{P}}
\newcommand{\cT}{\mathcal{T}}
\newcommand{\cX}{\mathcal{X}}
\newcommand{\cY}{\mathcal{Y}}

\newcommand{\sA}{\mathscr{A}}
\newcommand{\sB}{\mathscr{B}}
\newcommand{\sC}{\mathscr{C}}
\newcommand{\sE}{\mathscr{E}}
\newcommand{\sF}{\mathscr{F}}
\newcommand{\sG}{\mathscr{G}}
\newcommand{\sH}{\mathscr{H}}
\newcommand{\sL}{\mathscr{L}}
\newcommand{\sS}{\mathscr{S}}
\newcommand{\sT}{\mathscr{T}}
\newcommand{\sX}{\mathscr{X}}
\newcommand{\sY}{\mathscr{Y}}
\newcommand{\sZ}{\mathscr{Z}}

\newcommand{\tS}{\tilde{S}}
% Debug
\newcommand{\todo}[1]{\begin{color}{blue}{{\bf~[TODO:~#1]}}\end{color}}

% a few handy macros

\renewcommand{\le}{\leqslant}
\renewcommand{\ge}{\geqslant}
\newcommand\matlab{{\sc matlab}}
\newcommand{\goto}{\rightarrow}
\newcommand{\bigo}{{\mathcal O}}
%\newcommand{\half}{\frac{1}{2}}
%\newcommand\implies{\quad\Longrightarrow\quad}
\newcommand\reals{{{\rm l} \kern -.15em {\rm R} }}
\newcommand\complex{{\raisebox{.043ex}{\rule{0.07em}{1.56ex}} \hskip -.35em {\rm C}}}


% macros for matrices/vectors:

% matrix environment for vectors or matrices where elements are centered
\newenvironment{mat}{\left[\begin{array}{ccccccccccccccc}}{\end{array}\right]}
\newcommand\bcm{\begin{mat}}
\newcommand\ecm{\end{mat}}

% matrix environment for vectors or matrices where elements are right justifvied
\newenvironment{rmat}{\left[\begin{array}{rrrrrrrrrrrrr}}{\end{array}\right]}
\newcommand\brm{\begin{rmat}}
\newcommand\erm{\end{rmat}}

% for left brace and a set of choices
%\newenvironment{choices}{\left\{ \begin{array}{ll}}{\end{array}\right.}
\newcommand\when{&\text{if~}}
\newcommand\otherwise{&\text{otherwise}}
% sample usage:
%\delta_{ij} = \begin{choices} 1 \when i=j, \\ 0 \otherwise \end{choices}


% for labeling and referencing equations:
\newcommand{\eql}{\begin{equation}\label}
\newcommand{\eqn}[1]{(\ref{#1})}
% can then do
%\eql{eqnlabel}
%...
%\end{equation}
% and refer to it as equation \eqn{eqnlabel}.
% some useful macros for finite difference methods:
\newcommand\unp{U^{n+1}}
\newcommand\unm{U^{n-1}}
% for chemical reactions:
\newcommand{\react}[1]{\stackrel{K_{#1}}{\rightarrow}}
\newcommand{\reactb}[2]{\stackrel{K_{#1}}{~\stackrel{\rightleftharpoons}
 {\scriptstyle K_{#2}}}~}
\makeatletter
\def\th@plain{%
\thm@notefont{}% same as heading font
\itshape % body font
}
\def\th@definition{%
\thm@notefont{}% same as heading font
\normalfont % body font
}
\makeatother
\date{}


\title{Lecture-25: Poisson processes: Conditional distribution}
\author{}

\begin{document}
\maketitle


\section{Joint conditional distribution of points in a finite window}
Let $\sX =\R^d$ be a $d$-dimensional Euclidean space, 
and $S:\Omega\to\sX^\N$ be a Poisson point process with intensity measure $\Lambda: \sB(\sX) \to \R_+$ and associated counting process $N:\Omega\to\Z_+^{\sB(\sX)}$.
\begin{prop} 
Let $k \in \N$ be any positive integer. 
Consider a Poisson point process $S: \Omega \to \sX^\N$ with intensity measure $\Lambda:\sB(\sX)\to\R_+$,  
a finite partition $A\in \sB(\sX)^k$ that partitions a bounded set $B \in \sB(\sX)$, 
and a vector $n \in \Z_+^k$ that partitions a non negative integer $m \in \Z_+$. 
Then, 
\EQN{
\label{eqn:multi}
P(\set{N(A_1) = n_1, \dots, N(A_k) = n_k}\given\set{N(B) = m}) 
= 
\binom{m}{n_1, \dots, n_k}\prod_{i=1}^k\left(\frac{\Lambda(A_i)}{\Lambda(B)}\right)^{n_i}.
}
\end{prop}
\begin{proof} 
From the definition of conditional probability and the fact that $\cap_{i=1}^k\set{N(A_i) = n_i} \subseteq \set{N(B) =m}$, we can write the conditional probability on LHS as the ratio 
\begin{equation*}
%P(\set{N(A_1) = n_1, \dots, N(A_k) = n_k}\given\set{N(A) = n}) = 
\frac{P\set{N(A_1) = n_1, \dots, N(A_k) = n_k, N(B) = m}}{P\set{N(B) = m}} = 
\frac{P\set{N(A_1) = n_1, \dots, N(A_k) = n_k}}{P\set{N(B) = m}}.
\end{equation*}
From the complete independence property and Poisson marginals for the joint distribution of $(N(A_1), \dots, N(A_k))$ for the partition $A \in \sB(\sX)^k$, 
and the fact that the intensity measure sums over disjoint sets, i.e. $\Lambda(B) = \sum_{i=1}^k\Lambda(A_i)$, we can rewrite the RHS of the above equation as 
\begin{equation*}
\frac{P\set{N(A_1) = n_1, \dots, N(A_k) = n_k}}{P\set{N(B) = m}}
= \frac{\prod_{i=1}^ke^{-\Lambda(A_i)}\frac{\Lambda(A_i)^{n_i}}{n_i!}}{e^{-\Lambda(B)}\frac{\Lambda(B)^m}{m!}}
= \binom{m}{n_1, \dots, n_k}\prod_{i=1}^k\left(\frac{\Lambda(A_i)}{\Lambda(B)}\right)^{n_i}.
\end{equation*}
\end{proof}
\begin{rem}
Consider a Poisson point process $S:\Omega\to\sX^\N$ with intensity measure $\Lambda: \sB(\sX) \to \R_+$ and counting process $N:\Omega\to\Z_+^{\sB(\sX)}$. 
Let $A \in \sB(\sX)^k$ be a partition for bounded set $B \in \sB(\sX)$.  
\begin{enumerate}[i\_]
\item Defining $p_i \triangleq \frac{\Lambda(A_i)}{\Lambda(B)}$, we see that $(p_1, \dots, p_k) \in \cM([k])$ is a probability distribution. % over partitions $(A_1, \dots, A_k)$. 
We also observe that 
\EQ{
p_i 
= P(\set{N(A_i)=1}\given\set{N(B) = 1}) 
= P(\set{\abs{S \cap A_i} = 1}\given\set{\abs{S \cap B} = 1}).
}
When $N(B)=1$, we can call the point of $S$ in $B$ as $S_1$ without any loss of generality. 
That is, if we call $\set{S_1} = S \cap B$, then we have 
\EQ{
p_i = P(\set{S_1 \in A_i}\given\set{S_1 \in B}).
}
Similarly, when $N(B) = n_i$, 
we call the points of $S$ in $B$ as $S_1, \dots, S_{n_i}$ and denote $S\cap B = \set{S_1, \dots, S_{n_i}}$. 
For this case,  we observe
\eq{
&P(\set{N(A_i)=n_i}\given\set{N(B) = n_i}) 
= P(\set{\abs{S \cap A_i} = n_i}\given\set{\abs{S \cap B} = n_i}) = p_i^{n_i}\\
&= P(\cap_{j=1}^{n_i}\set{S_j \in A_i}\given \set{\set{S_1, \dots, S_{n_i}} = S\cap B}) 
= \prod_{j=1}^{n_i}P(\set{S_j \in A_i}\given \set{S_j \in B}). 
}

\item We can rewrite the Equation~\eqref{eqn:multi} as a multinomial distribution, where 
\EQ{
 P(\set{N(A_1) = n_1, \dots, N(A_k) = n_k}\given\set{N(B) = m}) = {m \choose n_1, \dots, n_k} p_1^{n_1} \dots p_k^{n_k}.
} 

\item For any finite set $F\subseteq \N$ of size $m \in \Z_+$ and $n\in\Z_+^k$ a partition of $m$, 
we define $\cP_k(F, n)$ to be the collection of all $k$-partitions $E \in \cP(\N)^k$ of $F$ such that $\abs{E_i} = n_i$ for $i \in [k]$. 
Then, the multinomial coefficient accounts for number of partitions of $m$ points into sets with $n_1,  \dots, n_k$ points. 
That is, 
\EQ{ 
{m \choose n_1, \dots, n_k} = \abs{\cP_k([m], n)}.
}

\item %Defining $E_i \triangleq S \cap A_i$ and $E = S \cap A$, 
Recall that the event $\set{N(A_i) = n_i} = \set{\abs{S \cap A_i} = n_i}$. 
Hence, we can write 
\eq{
P(\cap_{i=1}^k\set{\abs{S \cap A_i} = n_i}\given \set{\abs{S \cap B} = m}) 
&= {m \choose n_1, \dots, n_k} p_1^{n_1} \dots p_k^{n_k}\\
&= \sum_{E \in \cP_k(S \cap B, n)}\prod_{i=1}^k\prod_{S_j\in E_i}P(\set{S_j \in A_i}\given \set{S_j \in B}).
}

\item 
When $N(B) = m$, we denote $S\cap B$ by $F = \set{S_1, \dots, S_m}$ without any loss of generality. 
We further observe that when $N(A_i) = n_i$ for all $i \in [k]$, 
then $(S \cap A_1, \dots, S \cap A_k) \in \cP_k(S\cap B,n)$. 
Therefore, we can re-write the event 
\EQ{
\cap_{i=1}^k\set{N(A_i) = n_i} 
= \cap_{i=1}^k\set{\abs{S \cap A_i} = n_i} 
= \cup_{E \in \cP_k(S\cap B, n)}(\cap_{i=1}^k\set{S \cap A_i = E_i}).
} 
That is, we can write the conditional probability conditioned on $S\cap B = F$, as 
\eq{
P(\cap_{i=1}^k\set{N(A_i) = n_i} \given \set{N(B) = m}) 
&= \sum_{E \in \cP_k(F,n)}P(\cap_{i=1}^k\set{S \cap A_i = E_i}\given\set{S\cap B = F})\\
&= \sum_{E \in \cP_k(F, n_1, \dots, n_k)}P(\cap_{i=1}^k\cap_{S_j \in E_i}\set{S_j \in A_i}\given\set{S\cap B = F}).
}

\item   
Equating the RHS of the above equation term-wise, we obtain that conditioned on each of these points falling inside the window $B$, 
the conditional probability of each point falling in partition $A_i$ is independent of all other points and given by $p_i$. 
That is, we have 
\EQ{
 P(\cap_{i=1}^k\cap_{S_j \in E_i}\set{S_j \in A_i}\given\set{S\cap B = F}) 
 = \prod_{i=1}^k\prod_{S_j\in E_i}P(\set{S_j \in A_i}\given \set{S_j \in B}) = \prod_{i=1}^kp_i^{n_i} =  \prod_{i=1}^k\left(\frac{\Lambda(A_i)}{\Lambda(B)}\right)^{n_i} .
} 
It means that given $m$ points in the window $B$, the location of these points are independently and identically distributed in $B$ according to the distribution $\frac{\Lambda(\cdot)}{\Lambda(B)}$. 

%\item We can also see this by observing that the event $\set{N(A_i) = n_i} = \set{\sum_{S_j \in S}\indicator{S_j \in A_i} = n_i}$. 
%Hence, 
%\EQ{
%P(\set{N(A_1) = n_1, \dots, N(A_k) = n_k}\given\set{N(A) = n})  
%= P(\{\sum_{S_j \in S}\indicator{S_j \in A_i} = n_i, i \in [k]\}\given \set{N(A) = n}).
%}
%\item Let $A \in \sB$, its partitioning sets $A_1, \dots, A_k \in \sB$, 
%the numbers $N(A_i) = n_i = 1$ for all $i \in [k]$,  
%and a bounded set $W \in \sB$ such that $A \subseteq W$ and $N(W) = k$. 
%Recall that $\set{N(A_i) = 1} = \set{S_j \in A_i \text{ for exactly one }S_j \in S}$. 
%Given that $\set{N(A) =k}$, we can write the following set equality 
%\EQ{
%\set{N(A_1) = 1, \dots, N(A_k) = 1, N(W \setminus A) = 0} = \cup_{\sigma: [k] \to [k] \text{ permutation }}\set{S_{\sigma(1)} \in A_1, \dots, S_{\sigma(k)} \in A_k}.
%}
%Therefore, it follows that for any bounded disjoint subsets $A_1, \dots, A_k \in \sB$ such that $A = \cup_{i=1}^k A_i \subseteq W$
%\EQ{
%P(\set{S_1 \in A_1, \dots, S_k \in A_k} \given \set{S_1, \dots, S_k \in W}) = \prod_{i=1}^k\frac{\Lambda(A_i)}{\Lambda(W)} = \prod_{i=1}^kP(\set{S_i \in A_i}\given \set{S_i \in W}).
%}
\item If the Poisson process is homogeneous, the distribution is uniform over the window $B$. 
\item For a Poisson process with intensity measure $\Lambda$ and any bounded set $A \in \sB$, 
the number of points $N(A)$ in the set $A$ is a Poisson random variable with parameter $\Lambda(A)$. 
Given the number of points $N(A)$, the location of all the points in $S\cap A$ are \iid with density $\frac{\lambda(x)}{\Lambda(A)}$ for all $x \in A$. 
\end{enumerate}
\end{rem}
\begin{shaded}
\begin{rem}[Simulating a homogeneous Poisson point process] 
If we are interested in simulating a two dimensional homogeneous Poisson point process with density $\lambda$ in a uniform square $A = [0,1]\times[0,1]$. 
Then, we first generate the random variable $N(A):\Omega\to\Z_+$ that takes value $n$ with probability $e^{-\lambda}\frac{\lambda^n}{n!}$. 
Next, for each of the $N(A)=n$ points, we generate the location $(X_i, Y_i) \in \R^2$ uniformly at random. 
That is, $X:\Omega\to[0,1]^n$ and $Y:\Omega\to[0,1]^n$ are independent \iid uniform sequences.
\end{rem}
\end{shaded}
\begin{cor} 
For a homogeneous Poisson point process on the half-line with ordered set of points $\tilde{S} = (S_{(n)} \in \R_+: n \in \N)$, 
we can write the conditional density of ordered points $(S_{(1)}, \dots, S_{(k)})$ given the event $\set{N_t = k}$ as the ordered statistics of \iid uniformly distributed random variables. 
Specifically, we have 
\EQ{
f_{S_{(1)}, \dots, S_{(k)}\given N_t = k}(t_1, \dots, t_k) = k!\frac{\SetIn{0 < t_1 \le \dots \le t_k \le t}}{t^k}. 
}
\end{cor}
\begin{proof}
Given $\set{N_t = k}$, we can denote the points of the Poisson process in $(0,t]$ by $S_1, \dots, S_k$. 
From the above remark, we know that $S_1, \dots, S_k$ are \iid uniform in $(0, t]$, 
conditioned on the number of points $N_t = k$. 
Hence, we can write 
\EQ{
F_{S_1, \dots, S_k\given N_t = k}(t_1, \dots, t_k) 
= P(\cap_{i=1}^k\set{S_i \in (0, t_i]}\given \set{N_t=k}) 
= \prod_{i=1}^kP(\set{S_i \in (0, t_i]}\given \set{S_i \in (0, t]}) 
= \prod_{i=1}^k\frac{t_i}{t}\SetIn{0 < t_i \le t}.
}  
Therefore, for $0 < t_1 < \dots < t_k < 1$ and $(h_1, \dots, h_k)$ sufficiently small, we have 
\EQ{
P\Big(\cap_{i=1}^k\set{S_i \in (t_i, t_i+h_i]}\Big) = \prod_{i=1}^k\frac{h_i}{t}.
}
Since $(S_1, \dots, S_k)$ are conditionally independent given $S\cap A = \set{S_1, \dots, S_k}$, 
it follows that any permutation $\sigma: [k] \to [k]$, the conditional joint distribution of $(S_{\sigma(1)}, \dots, S_{\sigma(k)})$ is identical to that of $(S_1, \dots, S_k)$ given $S\cap A = \set{S_1, \dots, S_k}$. 
Further, we observe that the order statistics of $(S_{\sigma(1)}, \dots, S_{\sigma(k)})$ is identical to that of $(S_{1}, \dots, S_{k})$.  
Therefore, we can write the following equality for the events
%\EQ{
%\set{0 < S_{(1)} \le \dots \le S_{(k)} \le t_k \le t} 
%= \cup_{\sigma: [k] \to [k] \text{ permutation}}\cap_{i=1}^k\set{S_{\sigma(i)}\le t_i}. 
%}
\EQ{
\cap_{i=1}^k\set{S_{(i)}\in (t_i, t_i+h_i]} 
= \cup_{\sigma: [k] \to [k] \text{ permutation}}\cap_{i=1}^k\set{S_{\sigma(i)}\in (t_i, t_i+h_i]}. 
}
The result follows since the number of permutations $\sigma: [k] \to [k]$ is $k!$. 
%\EQ{
 %P(\cap_{i=1}^k\set{S_{(i)} \in (0, t_i]}\given \set{N_t=k})  
%= \sum_{\sigma: [k] \to [k]}P(\cap_{i=1}^k\set{S_{\sigma(i)} \in (0, t_i]}\given \set{N_t=k}).
%}
\end{proof}
%\begin{cor} 
%Let $(N_t: t \in \R_+)$ be the counting process on the half-line with ordered set of points $(S_{(n)}: n \in \N)$ from a homogeneous Poisson point process.   
%Let $N_t = k$, instants $0 = t_0 < t_1 < \dots < t_{k-1} < t_k \le t$, sufficiently small durations $0 < h_i < t_i -t_{i-1}$, and intervals $I_i = (t_i-h, t_i]$, 
%then 
%\EQ{
%P(\set{N(I_1) = 1, \dots, N(I_k)=1}\given \set{N_t = k}) = k! \frac{h_1}{t}\dots \frac{h_k}{t}. 
%}
%\end{cor}
%\begin{rem} 
%Let $I_i = (t_i - h, t_i]$ for sufficiently small $h \in \R_+$. 
%We notice that the event $\set{N(I_1) = 1, \dots, N(I_k) =1} = \set{S_{(1)} \in I_1, \dots, S_{(k)} \in I_k}$.  
%Taking limit as $h_i \to 0$, we see that the conditional density of ordered sequence $(S_{(1)}, \dots, S_{(k)})$ given $N_t = k$ is 
%given by
%\EQ{
%f_{S_{(1)}, \dots, S_{(k)}\given N_t = k}(t_1, \dots, t_k) = \frac{k!}{t^k}\indicator{0 < t_1 < \dots < t_k \le t}. 
%}
%There are $k!$ sequences $(S_1, \dots, S_k)$ that give the same ordered sequence $(S_{(1)}, \dots, S_{(k)})$. 
%It follows that 
%\EQ{
%%P(\set{N(t_1) = 1, \dots, N(t_k) - N(t_{k-1})=1}\given \set{N_t = k}) = \frac{1}{k!}P(\set{S_1 \le t_1, \dots, S_k \le t_k}\given\set{N_t = k}) = \prod_{i=1}^k\frac{t_i}{t}. 
%f_{S_{1}, \dots, S_{k}\given N_t = k}(t_1, \dots, t_k) = \prod_{i=1}^k\frac{\indicator{t_i \le t}}{t}. 
%}
%That is conditioned on $n$ points in the interval $(0, t]$, the un-ordered arrival points are all uniformly distributed in $(0, t]$, 
%and the ordered points are ordered statistics of these uniformly distributed points.  
%\end{rem} 


%\section{Poisson Point Processes} 
%A simple point process $S = \{S_i \in \R^d: i \in \N\}$ is a random countable collection of distinct points in the $d$-dimensional Euclidean space $\R^d$. 
%Let $\sB$ be the collection of Borel measurable sets in $\R^d$. 
%Then, the associated counting process $(N(A): A \in \sB)$ is defined as $N(A) = \sum_{i \in \N}1_{\{S_i \in A\}}$. 
%For a sample realization of a simple point process $S$, we observe that $dN(x) = 0$ for all $x \notin S$. 
%Hence, for any function $f : \R^d \to \R$, we have
%\EQ{
%\int_{x \in A}f(x)dN(x) = \sum_{i \in \N}f(S_i)1_{\{S_i \in A\}}.
%}
%%A simple point process can also be thought as a random process $(S(A): A \in \sF)$ indexed by all Boreal measurable sets of $\R^d$. 
%The finite dimensional distribution of $N$ is joint distribution of $(N(A_1), \dots, N(A_k))$ for some finite $k \in \N$ and bounded sets $A_1, \dots, A_k \in \sF$. 
%A counting process $N$ has the \textbf{completely independence property},  
%if for any collection of finite disjoint and bounded sets $A_1, \dots, A_k \in \sF$, 
%\EQ{
%P\bigcap_{i=1}^k\{N(A_i) = n_i\} = \prod_{i=1}^nP\{N(A_i) = n_i\}.
%}
%
%A \textbf{Poisson point process} of intensity measure $\Lambda$ is defined in terms of the finite-dimensional distribution of the associated counting process, as  
%\EQ{
%P\bigcap_{i=1}^k\{N(A_i) = n_i\} = \prod_{i=1}^k\left(e^{-\Lambda(A_i)}\frac{\Lambda(A_i)^n_i}{n_i!}\right),
%}
%for some finite $k \in \N$, mutually disjoint and bounded sets $A_1, \dots, A_k \in \sF$ and $\Lambda(A) \triangleq \int_{x \in A}d\Lambda(x)$. 
%
%\begin{thm} 
%Following are equivalent for a simple counting process $N$. 
%\begin{enumerate}[i\_]
%\item Process $N$ is Poisson with intensity measure $\Lambda$. 
%\item Process $N$ has the completely independence property, and $\E N(A) = \Lambda(A)$. 
%\item The intensity measure $\Lambda$ is bounded for bounded $A \in \sF$,  and $N(A)$ is a Poisson  with parameter $\Lambda(A)$. 
%\end{enumerate}
%\end{thm}
%
%\begin{prop} 
%For a finite $k \in \N$, disjoint bounded sets $A_1, \dots, A_k \in \sF$, numbers $n_1, \dots, n_k \in \Z_+$, let $A = \cup_{i=1}^kA_i$ and $n = \sum_{i=1}^kn_i$. 
%Then for a Poisson point process, we have 
%\EQ{
%P(N(A_1) = n_1, \dots, N(A_k) = n_k | N(A) = n) = \frac{n!}{n_1!\dots n_k!}\prod_{i=1}^n\left(\frac{\Lambda(A_i)}{\Lambda(A)}\right)^{n_i}. 
%}
%\end{prop}
%\begin{rem} 
%Above proposition implies that $n$ points in a bounded set $A$ are \emph{iid} distributed in this set with density function $\frac{\Lambda'(x)}{\Lambda(A)}$ for $x \in A$. 
%\end{rem} 


%\appendix 
%\section{Functions with semigroup property}
%\begin{lem} 
%\label{lem:semigroup}
%A unique non-negative right continuous function $f: \R_+ \to \R_+$ satisfying the semigroup property
%\EQ{
% f(t+s) = f(t)f(s), \text{ for all } t,s \in \R_+
% }
% is $f(t) = e^{\theta t}$, where $\theta = \log f(1)$.
%\end{lem}
%\begin{proof}
%Clearly, we have $f(0) = f^2(0)$. Since $f$ is non-negative, it means $f(0) = 1$. 
%By definition of $\theta$ and induction for $m,n \in \Z^+$, we see that 
%\begin{xalignat*}{3}
%&f(m) = f(1)^m = e^{\theta m},&& e^{\theta} = f(1) = f(1/n)^n.
% \end{xalignat*}
%Let $q \in \Q$, then it can be written as $m/n, n \neq 0$ for some $m,n \in \Z^+$. 
%Hence, it is clear that for all $q \in \Q^+$, we have $f(q) = e^{\theta q}$.
%either unity or zero. Note, that $f$ is a right continuous function and is non-negative. 
%Now, we can show that $f$ is exponential for any real positive $t$ by taking a sequence of rational numbers $(q_n: n \in \N)$ decreasing to $t$. 
%From right continuity of $f$, we obtain 
%\EQ{
%f(t) = \lim_{q_n \downarrow t} f(q_n) =   \lim_{q_n \downarrow t} e^{\theta q_{n}}= e^{\theta t}.
%}
%\end{proof}

\end{document}